% =====================================================================================================
% Chương 4, PHẦN A (mở chương, mục 4.1–4.5) — bản sửa v4 (2026-10-12), BẢN DÙNG ĐỂ GHÉP.
% Theo ba bản nhận xét R25 (Codex, kiểm chứng sự kiện, người đọc và giảng viên); lý do từng thay đổi:
% reviews/CH4A-REVISION-LOG.md; ghi chú: drafts/08-chuong4-phanA-v2-NOTES.md.
% Bản 1 (08-chuong4-phanA.tex), v2 (lượt sửa đầu) và v3 (lượt rút gọn thứ nhất) giữ nguyên; v4 chỉ rút gọn câu
% chữ để vừa ngân sách trang; mọi chỗ cắt so với v2 ghi ở NOTES. Phần B (mục 4.6–4.11) nối ngay sau tệp này.
% Mỗi dòng "% EV:" trỏ tới EVIDENCE-REGISTER.md (A01…K52), EVIDENCE-REGISTER-ADDENDUM-1.md (N-xx, L1a-xx,
% mục 0.x), EVIDENCE-REGISTER-ADDENDUM-2-CH3.md (X3-xx) hoặc EVIDENCE-REGISTER-ADDENDUM-4-CH4A.md (C4A-xx).
% Không dùng gói mới: hình vẽ bằng TikZ thường (thư viện shapes.geometric đã nạp trong main.tex).
% =====================================================================================================

\chapter{Thực nghiệm và đánh giá}
\label{chuong:thucnghiem}

Chương này đo luật C được gì và mất gì trên mô phỏng FleetPy, so với điều phối không trần, luật hạn chót và cách đệm
lời hứa, tức cộng thêm một khoảng cố định vào mọi lời hứa. Chương cũng xét lời hứa theo dự báo và những chỗ dự báo
chưa đủ. Các mục đi theo ba phần của độ lệch giờ trả so với lời hứa đầu (mục~\ref{muc:so}). Phần tính cho điều phối,
gắn với việc ghép thêm khách, ở mục~\ref{muc:c4-ketqua}--\ref{muc:c4-vaitro}; phần được hồ sơ tắc đường lịch sử giải
thích, cùng cách đệm lời hứa, ở mục~\ref{muc:c4-dubao} và~\ref{muc:c4-dem}; phần còn lại ở mục~\ref{muc:c4-tuyet}.
Bảng~\ref{bang:c4-lotrinh} nối ba câu hỏi nghiên cứu ở mục~\ref{muc:c1-cauhoi} với các mục trả lời chúng.
% EV: — (đoạn dẫn, không nêu số); ba phần theo eq:so của chương 3 (X3-12); "đệm" theo GLOSSARY-VI mục B

\begin{table}[htbp]
\centering
\caption{Ba câu hỏi nghiên cứu và các mục của Chương~\ref{chuong:thucnghiem} trả lời chúng}
\label{bang:c4-lotrinh}
\small
\begin{tabular}{@{}p{6.3cm}p{2.4cm}p{6.3cm}@{}}
\toprule
\textbf{Câu hỏi} & \textbf{Mục} & \textbf{Đại lượng được đo}\\
\midrule
CH1. Phần dời do điều phối được đo thế nào, và trần trên phần đó giữ được gì, không giữ được gì?
& \ref{muc:c4-dieuphoi}, \ref{muc:c4-congdon}, \ref{muc:c4-hople}
& Số khách bị dời giờ trả vì ghép thêm khách, độ dời lớn nhất; yêu cầu bị từ chối có phương án bị trần chặn;
phép cộng các lần dời qua cả chuyến; các bộ đếm kiểm tra\\
CH2. Giữ trần đổi số lượt phục vụ, số lượt đúng giờ và độ dài lời hứa thế nào so với điều phối không
trần, luật hạn chót và cách đệm lời hứa?
& \ref{muc:c4-ketqua}, \ref{muc:c4-chontran}, \ref{muc:c4-vaitro}, \ref{muc:c4-dem}
& Lượt phục vụ; lượt trễ và không trễ so với lời hứa đầu; quyết định bị ép; các mức trần; thứ tự tiêu chí
của bộ chọn; độ dài thêm của lời hứa\\
CH3. Lời hứa theo hồ sơ tắc đường lịch sử đổi số lượt đúng giờ thế nào, và phần nào còn chưa được
giải thích?
& \ref{muc:c4-dubao}, \ref{muc:c4-tuyet}
& Lượt trễ khi có và không có dự báo; độ dài thêm của lời hứa; ba phần của độ lệch trên ngày thường và chiều
tuyết\\
\bottomrule
\end{tabular}
\end{table}
% EV: — (lộ trình, không có khẳng định; câu hỏi theo OUTLINE-V2 cuối §1.2)

% -----------------------------------------------------------------------------------------------------
\section{Thiết lập thực nghiệm}
\label{muc:c4-thietlap}

\subsection{Kịch bản thử}
\label{muc:c4-kichban}

FleetPy 1.0.2 mô phỏng xe và khách trên bộ dữ liệu chuẩn Manhattan; ở mỗi lượt quyết định, bộ chuyển đổi gửi tình
trạng hiện tại cho bộ chạy RideBound và nhận lại quyết định (mục~\ref{muc:kientruc}). Nhu cầu là các chuyến taxi New
York tuần 11--18/11/2018 trong bộ dữ liệu này~\cite{engelhardt2026}. Một kịch bản gồm một ngày, một khung giờ hai
tiếng và một trong bốn mẫu nhu cầu của ngày đó, mỗi mẫu gồm 10\% số chuyến của ngày: 108 yêu cầu trong khung giờ, phục
vụ bởi 8 xe 4 chỗ (Bảng~\ref{bang:c4-thamso}).
% EV: B01; X3-15; X3-21 (FC:395-429); B02 (FIX: 2018-11-14_sample_10_3.csv); J18 (mẫu 10% của một ngày); B08 (mẫu r1–r4); B03

Tắc đường là một hệ số chung cho cả thành phố, đổi theo khung 15 phút và lấy từ dữ liệu chuyến taxi, không phải tốc độ
từng đoạn đường; dưới đây gọi là hệ số tắc đường đo được. Thế giới mô phỏng không có nhiễu ngẫu nhiên, nên mỗi cặp
kịch bản và cấu hình chạy một lần và cho một kết quả. Ba kịch bản cùng ngày, cùng mẫu có chung 3 đến 20 trong 108 yêu
cầu, nên các kịch bản không độc lập với nhau.
% EV: B06 (sigma = shockRate = 0); K49; B11 (T38:90, 93); I05

\begin{table}[htbp]
\centering
\caption{Tham số của một kịch bản thử. Các giá trị là cấu hình của thực nghiệm, không phải chuẩn ngành. Tập tắc
đường cố định (mục~\ref{muc:c4-taptu}) khác ở dòng tắc đường và chạy bằng bộ chạy cũ, có giới hạn bộ sinh và bộ
giải chưa được kiểm lại.}
\label{bang:c4-thamso}
\small
\begin{tabular}{@{}p{4.0cm}p{11.0cm}@{}}
\toprule
\textbf{Tham số} & \textbf{Giá trị}\\
\midrule
Khung giờ & 7:00--9:00, 8:00--10:00 hoặc 17:00--19:00, gọi là khung 7, 8 và 17 giờ; riêng chiều tuyết
là 14:00--16:00 và 15:00--17:00\\
Yêu cầu & 108 mỗi kịch bản, lấy trong khung giờ từ một mẫu 10\% số chuyến taxi\\
Xe & 8 xe, mỗi xe 4 chỗ\\
Đồ thị & 96 nút là các điểm đi, điểm đến lấy từ dữ liệu Manhattan; mỗi cặp nút có sẵn một thời gian đi. Đây là
đồ thị rút gọn, không phải mạng đường phố\\
Ràng buộc dịch vụ & Đón trong vòng 600 giây kể từ lúc gọi; ngồi xe không quá 1,5 lần thời gian đi
thẳng (Bảng~\ref{bang:rangbuoc})\\
Thời gian dừng đón, trả & 0 giây: xe không mất thời gian khi khách lên, xuống (kiểm trên 3 lần chạy mẫu)\\
Tắc đường & Một hệ số chung cho cả thành phố theo khung 15 phút; tốc độ gốc của mọi kịch bản lấy theo hệ số lúc
00:00\\
Bộ sinh phương án & Tối đa 100 phương án mỗi xe; mỗi lượt quyết định xét tối đa 2 yêu cầu mới\\
Bộ giải & OR-Tools 9.15, hạt giống ngẫu nhiên cố định bằng 7\\
\bottomrule
\end{tabular}
\end{table}
% EV (khung giờ, yêu cầu): B08; B20; B02; C4A-06 (chiều tuyết cũng 108 yêu cầu mỗi kịch bản, INV:14-15,17-22)
% EV (xe, đồ thị, ràng buộc dịch vụ): B03; K48 (PREF:22-26: mỗi cặp điểm đi–đến là một cạnh trực tiếp)
% EV (dừng, tắc đường): B05 (chỉ 3 lần chạy mẫu, SERV:2-4); B06; B07 (tốc độ nền lấy lúc 00:00)
% EV (bộ sinh, bộ giải): A24; X3-17; X3-32 (giới hạn 2 yêu cầu áp cho cả lô của lượt, GEN:54-96); C4A-18 (tập tắc đường cố định: bộ chạy khác, NEW:13; giới hạn chưa kiểm)

\subsection{Các tập kịch bản và cấu hình được so}
\label{muc:c4-taptu}

Chương này dùng bốn tập kịch bản (Bảng~\ref{bang:c4-taptu}). Từ các ngày 14--18/11/2018 dựng được 49 kịch bản, chia
theo cặp ngày và mẫu để ba khung giờ của một cặp nằm cùng một phía. 34 kịch bản làm tập kiểm tra chính; 15 kịch bản
dành cho một phép hiệu chỉnh không dùng trong khóa luận. Ngày 14--16/11 là ngày thường, 17 và 18/11 là cuối tuần.
% EV: B08 (SPLIT: chia theo nhóm ngày, mẫu); C17 (15 kịch bản hiệu chỉnh khoảng dự phòng γ); B09; B21 (11/11/2018 là Chủ nhật)

\begin{table}[htbp]
\centering
\caption{Bốn tập kịch bản dùng trong chương. U: điều phối không trần; C-$B$: luật C trần $B$ giây;
M$^{+}$-$B$: luật hạn chót $B$ giây; V-30: luật trần độ dời khách thấy 30 giây; FREE: ghép tự do
(Bảng~\ref{bang:luat}). Cột cuối là số lần chạy, mỗi lần một cặp kịch bản và cấu hình.}
\label{bang:c4-taptu}
\footnotesize
\setlength{\tabcolsep}{3pt}
\begin{tabular}{@{}p{2.5cm}p{1.9cm}p{3.1cm}p{3.0cm}p{2.6cm}p{1.9cm}@{}}
\toprule
\textbf{Tập} & \textbf{Quy mô} & \textbf{Ngày, khung giờ} & \textbf{Cấu hình dùng ở chương này} &
\textbf{Bộ chạy} & \textbf{Lần chạy}\\
\midrule
Tập kiểm tra chính & 34 kịch bản, 3.672 yêu cầu & 14--18/11/2018; khung 7, 8, 17 giờ có 9, 14, 11
kịch bản & U, C-30, C-120, M$^{+}$-30, M$^{+}$-60, V-30 & Bản đã lưu, có luật phục hồi không tệ hơn &
170 (năm cấu hình) và 34 (C-120)\\
% EV: B08; B09; B10 (170 = 34 × 5); L1a-01 (34 lần chạy C-120); A19; A28; X3-10 (R120:3-8)
Tập chỉnh trần buổi sáng & 8 kịch bản, 864 yêu cầu & 12--13/11/2018; khung 7 giờ & U; C và M$^{+}$ ở 30, 60,
90, 120, 150, 180, 240 và 300 giây & Như tập chính & 136 (trong 192 lần chạy đã lưu)\\
% EV: B15; B16; D01; D04; X3-10 (RW7:1-8; RW7X:6-8); X3-28; C4A-07 (17 cấu hình × 8 = 136; 88 + 104 = 192, RW7:33-34; RW7X:35-36)
Tập tắc đường cố định & 16 kịch bản, 1.728 yêu cầu & 12--13/11/2018; khung 8 và 17 giờ; hệ số không
đổi trong lần chạy & FREE, U, C-30, C-60, C-120 & Bản cũ, không có luật phục hồi & 16 mỗi cấu hình\\
% EV: B17; D08; D09; H13 (n = 16 mỗi cấu hình); LAD:2-6 (travelSnapshotVersion {1: 16}: một bảng giờ đi); C4A-04
Chiều tuyết 15/11/2018 & 8 kịch bản, 864 yêu cầu & 15/11/2018; khung 14 và 15 giờ & C-30 &
Chưa kiểm lại & 8\\
% EV: B20; N-26 (8 lần chạy C-30, 601 khách được phục vụ); C4A-06 (864 = 8 × 108, INV:14-15,17-22); X3-10 (bộ chạy chiều tuyết chưa kiểm)
\bottomrule
\end{tabular}
\end{table}

Bản đã lưu là bản bộ chạy biên dịch và lưu lại, dùng chung cho tập kiểm tra chính, tập buổi sáng và các lần chạy trần
120 giây; bản cũ chưa có luật phục hồi (mục~\ref{muc:doichung}, \ref{muc:kiemtra}). Số của chiều tuyết mang lưu ý bộ
chạy chưa được kiểm lại. Trừ ghép tự do, mọi cấu hình dùng cùng bộ sinh, bộ kiểm và bộ chọn với luật C; điều phối
không trần chỉ khác luật C trần 30 giây ở trần giờ trả. Ghép tự do không khóa gì và không qua bộ kiểm cam kết.
% EV: A28; X3-10 (RW7:1-8; RW7X:6-8; R120:3-8; bộ chạy chiều tuyết chưa kiểm); B17; NEW:13 (bộ chạy cũ không có commitmentRecovery); A08; H01; A17; K50; X3-38 (HO:109, b1-rolling-cost)

Lần chạy hỏng vì lỗi hạ tầng (một lần ở tập kiểm tra chính) được thay bằng bản chạy lại hoàn tất (Phụ
lục~\ref{phuluc:C}). Bản kiểm kê ngày 08/10/2026 ghi 1.991 bản ghi lần chạy; kết quả của chương chỉ lấy từ bốn tập
trên, các đợt chạy khác chỉ được nhắc khi nêu phạm vi kiểm tra (mục~\ref{muc:c4-congdon}).
% EV: C4A-17 (T38:12; ADD1:6: C-30-d20181114-s10-r3-w17-rerun1); ADDENDUM-1 mục 0.3; B19; C4A-20 (INV:1, bản kiểm kê ngày 08/10)

\subsection{Chỉ số}
\label{muc:c4-chiso}

Một lượt phục vụ là một yêu cầu, tức một khách, được nhận và chở tới nơi. Lượt trễ và lượt tới sớm được đếm trong nhóm
khách được phục vụ, so với lời hứa đầu (Bảng~\ref{bang:c4-chiso}). Số lượt không trễ bằng số lượt phục vụ trừ số lượt
trễ, tính trên toàn bộ yêu cầu, nên yêu cầu bị từ chối không được tính là đúng giờ. Ngưỡng chính là 60 giây; ngưỡng
120 giây được báo kèm vì có phép so đổi kết luận theo ngưỡng (mục~\ref{muc:c4-ketqua}).
% EV: C01; N-07 (khách được nhận = khách được phục vụ); C06; C08 (DIALPY:8-9: bị từ chối không tính là kịp; ở 120 s U hơn C-30 5 lượt)

Lời hứa đầu là giờ trả của kế hoạch lúc khách xác nhận đặt xe, tính với thời gian đi lúc đó. Tính từ lúc đặt xe, nó
dài trung bình 14,15 phút khi không trần và 14,12 phút dưới luật C trần 30 giây (gộp theo khách), nên luật C không hứa
xa hơn.
% EV: A13; A05; B14 (BURD:5,11: promised door-to-door 14,15 / 14,12 phút)

Tỉ lệ \emph{gộp theo khách} chia tổng số khách thỏa điều kiện cho tổng số khách được phục vụ; \emph{trung bình theo
kịch bản} lấy trung bình của tỉ lệ từng kịch bản, mỗi kịch bản tính như nhau. Ví dụ dựng tay: với 5 trong 10 khách ở
một kịch bản và 5 trong 90 ở kịch bản kia, gộp theo khách cho 10\%, trung bình theo kịch bản cho 27,8\%. Chênh lệch là
cấu hình được thử trừ đối chứng; số dương nghĩa là nhiều hơn.
% EV: C27 (hai cách gộp cho hai số khác nhau); K24; ví dụ dựng tay (10/100 = 10%; (50% + 5,56%)/2 = 27,78%)

Kết quả ``có dự báo'' là phép phát lại trên các lần chạy đã ghi (mục~\ref{muc:dubao}): xe nào chở ai giữ nguyên, chỉ
các giờ báo cho khách được tính lại theo dự báo, và các chỉ số dùng $D_0$. Dự báo chỉ có ở 21 kịch bản ngày thường của
tập kiểm tra chính; 13 kịch bản cuối tuần và tập buổi sáng giữ lời hứa không dự báo.
% EV: C09; F03 (chỉ vế phát lại); X3-11 (D = max(p, S) áp cho mọi lời hứa đón và trả đã công bố); F02; N-28; D02 (12–13/11 không có dự báo)

\begin{table}[htbp]
\centering
\caption{Các chỉ số dùng trong chương. $a$ là giờ trả thực tế, $p_0$ là lời hứa đầu không dự báo, $D_0$
là lời hứa đầu theo dự báo (không có dự báo thì $D_0 = p_0$); $c$ là một lần dời giờ trả do điều phối, $z$ là phần
khách thấy của một lần cập nhật lời hứa (mục~\ref{muc:bamoc}). ``Quá'' nghĩa là lớn hơn hẳn: trễ đúng 60 giây
không tính là trễ quá 60 giây.}
\label{bang:c4-chiso}
\footnotesize
\setlength{\tabcolsep}{3pt}
\begin{tabular}{@{}p{2.7cm}p{5.3cm}p{3.5cm}p{3.5cm}@{}}
\toprule
\textbf{Chỉ số} & \textbf{Định nghĩa} & \textbf{Mẫu số, cách gộp} & \textbf{Vì sao đo}\\
\midrule
Lượt phục vụ & Yêu cầu được nhận và chở tới nơi & Số đếm; chia cho số yêu cầu thì ra tỉ lệ phục vụ & Giá của
luật tính bằng chuyến\\
Trễ quá $T$ & Khách được phục vụ có $a - D_0 > T$; $T$ = 30, 60, 120, 300 hoặc 600 giây & Đếm trong khách
được phục vụ & Lời hứa bị vượt ở phía muộn\\
Không trễ quá $T$ & Lượt phục vụ trừ lượt trễ quá $T$ & Chia cho số yêu cầu; yêu cầu bị từ chối không
tính là đúng giờ & Phục vụ và đúng giờ trên một mẫu số\\
Tới sớm quá 60 giây & Khách được phục vụ có $a - D_0 < -60$ giây & Đếm trong khách được phục vụ & Lời
hứa quá dài cũng sai\\
Lời hứa dài thêm & $D_0 - p_0$ & Trung bình gộp theo khách được phục vụ, kể cả khách không có dự báo & Giá của
dự báo hay đệm\\
Khách đi chung & Có khách khác trên cùng xe trong khoảng từ lúc được đón tới lúc được trả & Gộp theo
khách được phục vụ & Mức ghép chuyến còn lại\\
Khách bị điều phối dời giờ trả & Có ít nhất một lần $c \neq 0$; độ dời của khách là $\sum|c|$ &
Gộp theo khách được phục vụ & Đại lượng luật C giới hạn\\
Quyết định bị ép & Lượt quyết định có ít nhất một xe bị ép: luật lời hứa đã loại cả phương án giữ nguyên kế hoạch
của xe, và xe đi theo phương án được giữ nhờ luật phục hồi, có ít nhất một khách được miễn trừ
(mục~\ref{muc:doichung}) & Trong mỗi kịch bản, chia cho số lượt quyết định đã được bộ mô phỏng xác nhận áp dụng;
rồi lấy trung bình theo kịch bản. Số gộp ghi riêng & Luật có tự giữ được nghĩa vụ không\\
Khách bị phá lời hứa & Khách mà luật đang dùng không giữ được điều nó hứa. Luật C: $\sum|c| > B$; luật hạn chót:
$a - p_0 > B$; luật trần độ dời khách thấy: $\sum|z| > B$ & Trong mỗi kịch bản, chia cho số khách được phục vụ;
rồi lấy trung bình theo kịch bản. Số gộp ghi riêng & Luật có giữ được điều nó hứa không\\
\bottomrule
\end{tabular}
\end{table}
% EV (phục vụ): C01; L1a-01 (2.387 hoàn tất + 1.285 bị từ chối = 3.672); N-07
% EV (trễ, không trễ, sớm): C06; C08 (DIALPY:8-9); N-04 (trễ > 60 s và sớm < −60 s rời nhau); C4A-01 (ngưỡng 30 s, cùng quy tắc sweep_analyze.py:107,128-132); C09; C10
% EV (dài thêm, đi chung, bị dời): F06; N-06 (gộp theo khách, ext_sum / served, ADD1:14,19,24); N-09 (định nghĩa đi chung); C20; N-07; A06; X3-05
% EV (bị ép): C4A-10 (RS:1012-1027; VALID:141-145; HVA:252-300; POL:443-463); X3-30; X3-09; C13 (trung bình theo kịch bản)
% EV (phá lời hứa): N-21 (định nghĩa của M⁺-30 và C-30; số gộp); A11; GLOSSARY-VI mục B "phá lời hứa"

\subsection{Mức trần 30 giây và tính chất thăm dò của kết quả}
\label{muc:c4-thamdo}

Trần chính $B$ = 30 giây là trần giờ trả của cấu hình có trần ở tập kịch bản tĩnh của đợt thử trước, nơi thời gian đi
không đổi (Phụ lục~\ref{phuluc:A}). Mức này được giữ nguyên cho tập kiểm tra chính, không chỉnh theo kết quả của tập
này. Nó không phải ngưỡng cảm nhận của khách hay chuẩn ngành, nên mỗi mức trần chỉ là một điểm trên đường đánh đổi
(mục~\ref{muc:c4-chontran}).
% EV: C4A-15 (TEX08:136-137: C1 = trần 30 giây cho tổng độ dời giờ trả; thay X3-18 đã bỏ); K29 (tập thời gian đi tĩnh); K33; B10 (năm cấu hình cố định trước lần chạy đầu, T38:3-6); D20; K52; K36

Các kết quả là thăm dò và dùng lại dữ liệu đã xem. Ngưỡng phân tích và cách chia tập được chọn sau khi đã xem tập buổi
sáng và các đợt chạy thử trước đó. Mức thử 120 giây được chọn sau khi đã xem kết quả của điều phối không trần và luật
C trần 30 giây trên chính tập kiểm tra chính, kể cả các phép sàng lọc ở 120 giây (mục~\ref{muc:c4-dieuphoi},
\ref{muc:c4-congdon}). Các con số vì vậy mô tả điều đã đo, không phải phép kiểm trên dữ liệu chưa từng xem.
% EV: I06; T38:8 (chọn sau khi xem Việc A, B, w07x); C4A-11 (BUD 07/10 = D14; CVP 10/10 = E06; C120R:10 chạy 10/10; T38:5 tập chính chạy 1–2/10); K42

% -----------------------------------------------------------------------------------------------------
\section{Kết quả chính khi chưa có dự báo}
\label{muc:c4-ketqua}

Trên 3.672 yêu cầu của tập kiểm tra chính, luật C với trần 30 giây không để phần dời giờ trả do điều phối của khách
nào vượt 30 giây: lớn nhất 29,2 giây, so với 590,1 giây (gần 10 phút) khi không trần. Không lượt nào của luật C trễ
quá 5 phút so với lời hứa đầu, so với 18 lượt khi không trần. Đổi lại, luật C phục vụ 2.338 lượt, ít hơn điều phối
không trần 96 lượt, và có ít hơn 242 khách đi chung. Số lượt không trễ quá một phút tăng 31, từ 2.231 lên 2.262, nhưng
số lượt không trễ quá hai phút giảm 5, từ 2.332 xuống 2.327.
% EV: C20; N-23; N-24; C07; C01; C02; N-20; C08

Mục~\ref{muc:c4-dem} cho thấy khi có dự báo, cộng thêm 30 giây vào mọi lời hứa của điều phối không trần cũng cho 2.328
lượt không trễ quá một phút, bằng luật C có dự báo. Vì vậy khóa luận không dùng chỉ số này làm lý do chọn luật C.
% EV: C10 (U + dự báo + 30 s: 2.328 = C-30 + dự báo: 2.328); K06; K51

Luật hạn chót chỉ mượn phần hạn chót của Schulz và Pfeiffer~\cite{schulz2026}: giờ trả mới không được muộn hơn lời hứa
đầu quá 30 hoặc 60 giây, thay vì 900 giây của họ. Nó chạy cùng bộ điều phối và luật phục hồi với luật C, nên đây không
phải phép so với hệ thống của hai tác giả.
% EV: J10; J11; A09; A19; H01; X3-28; X3-41

\begin{table}[htbp]
\centering
\caption{Kết quả trên tập kiểm tra chính với lời hứa đầu không dự báo (34 kịch bản, 3.672 yêu cầu mỗi
cấu hình). Số đếm gộp trên mọi kịch bản; trễ tính so với lời hứa đầu. Hai cột ``không trễ quá'' là số lượt phục vụ
không trễ quá ngưỡng; chia cho 3.672 thì ra tỉ lệ trên yêu cầu. ``Chưa đếm'': chưa có số đếm chính xác.}
\label{bang:c4-chinh}
\small
\begin{tabular}{@{}lrrrrr@{}}
\toprule
\textbf{Cấu hình} & \textbf{Phục vụ} & \textbf{Trễ quá 60 s} & \textbf{Không trễ quá 60 s} &
\textbf{Không trễ quá 120 s} & \textbf{Khách đi chung}\\
\midrule
U & 2.434 & 203 & 2.231 & 2.332 & 555\\
C-30 & 2.338 & 76 & 2.262 & 2.327 & 313\\
M$^{+}$-30 & 2.338 & 74 & 2.264 & 2.327 & 307\\
% EV (U, C-30, M⁺-30): C01; C06; C08; N-09 (555 / 313 / 307, số đếm chính xác)
M$^{+}$-60 & 2.347 & 79 & 2.268 & 2.336 & chưa đếm\\
V-30 & 2.336 & 74 & 2.262 & 2.325 & chưa đếm\\
% EV (M⁺-60, V-30): C01; C4A-02 (CK4A check: trễ quá 60 / 120 s 79 / 11 và 74 / 11; EXP:5-6 trễ quá 60 s 79 và 74; đi chung chỉ có % làm tròn ở EXP)
\bottomrule
\end{tabular}
\end{table}

Cách đọc bảng: dưới trần 30 giây, 2.338 yêu cầu được phục vụ; 76 lượt trong số đó tới muộn hơn lời hứa đầu quá 60
giây, nên 2.338 $-$ 76 = 2.262 lượt không trễ quá 60 giây.
% EV: C01; C06; C08

Cái giá trước hết là số lượt phục vụ: bớt 96 lượt nghĩa là tỉ lệ phục vụ giảm 2,6 điểm phần trăm (96 trên 3.672 yêu
cầu), tức ít hơn 3,9\% số lượt của điều phối không trần (96 trên 2.434). Mức giảm nằm ở khung 8 giờ (1.014 xuống 975)
và 17 giờ (904 xuống 852); khung 7 giờ chỉ giảm 5. Số khách đi chung giảm từ 555 xuống 313, tức từ 22,8\% xuống 13,4\%
số khách được phục vụ (gộp theo khách). Khóa luận không có bằng chứng về giá trị của các chuyến đi chung này đối với
khách hay hãng xe.
% EV: C02; C03; N-09; N-20 (lưu ý: D20, I07, I10)

Số lượt không trễ quá một phút vẫn tăng vì phục vụ bớt 96 lượt nhưng trễ bớt 127 lượt: 127 $-$ 96 = 31. Hai cấu hình
cũng phục vụ những khách khác nhau: 2.207 khách chung, 227 khách chỉ được điều phối không trần và 131 khách chỉ được
luật C phục vụ. Vì vậy 31 là chênh lệch hai tổng, không phải 31 khách cụ thể được đưa tới đúng giờ.
% EV: C06; C08; C04 (câu về chờ, ngồi xe của khách chung, C05, cắt vì ngân sách trang; xem NOTES)

So với luật hạn chót 30 giây, luật C phục vụ đúng bằng số lượt (2.338) nhưng không hoàn toàn cùng những khách (2.332
khách chung, mỗi bên 6 khách riêng), và có 76 so với 74 lượt trễ quá một phút. Luật hạn chót 60 giây và luật trần độ
dời khách thấy cũng gần như vậy (Bảng~\ref{bang:c4-chinh}). Các luật lời hứa khác nhau ở chỗ mỗi luật có giữ được điều
chính nó hứa hay không (Bảng~\ref{bang:c4-nghiavu}).
% EV: C01; C04; C06; K05; K07; C32; C4A-02

\begin{table}[htbp]
\centering
\caption{Nghĩa vụ của chính từng luật trên tập kiểm tra chính. Hai cột tỉ lệ là trung bình theo kịch bản (34
kịch bản); trong mỗi kịch bản, mẫu số lần lượt là số lượt quyết định được xác nhận và số khách được phục vụ. Số gộp
ghi riêng trong lời văn. Định nghĩa ``bị ép'' và ``bị phá lời hứa'' của từng luật ở Bảng~\ref{bang:c4-chiso}. Ở
hàng luật C, các số 0 là điều Nhận xét~\ref{nx:giunguyen} cho trước dưới giả thiết (N).}
\label{bang:c4-nghiavu}
\small
\begin{tabular}{@{}lrrr@{}}
\toprule
\textbf{Cấu hình} & \textbf{\begin{tabular}[b]{@{}r@{}}Quyết định\\ bị ép (\%)\end{tabular}} &
\textbf{\begin{tabular}[b]{@{}r@{}}Khách bị phá\\ lời hứa (\%)\end{tabular}} &
\textbf{\begin{tabular}[b]{@{}r@{}}Kịch bản có ít nhất\\ một khách bị phá\end{tabular}}\\
\midrule
C-30 & 0 & 0 & 0\\
M$^{+}$-30 & 41,3 & 14,5 & 28\\
M$^{+}$-60 & 22,5 & 4,2 & 22\\
V-30 & 53,8 & 21,6 & 34\\
\bottomrule
\end{tabular}
\end{table}
% EV: C13 (trung bình theo kịch bản; kịch bản có ≥ 1 phá hứa 0/28/22/34); N-21 (C-30 0/34 kịch bản, 0 quyết định bị ép); X3-30; A22

Cách đọc bảng: với luật hạn chót 30 giây, trung bình mỗi kịch bản có 41,3\% lượt quyết định bị ép và 14,5\% khách được
phục vụ tới muộn hơn lời hứa đầu quá 30 giây. Có 28 trong 34 kịch bản có ít nhất một khách như vậy.
% EV: C13; N-21

Luật C không có quyết định nào bị ép và không vượt trần của khách nào ở cả 34 kịch bản. Đây là điều Nhận
xét~\ref{nx:giunguyen} cho trước dưới giả thiết (N), không phải phát hiện thực nghiệm: phương án giữ nguyên kế hoạch
có phần dời do điều phối bằng 0 nên không bị luật C loại. Điều phối không trần cũng có 0 trong 20.760 quyết định bị
ép, và nhận xét không bảo đảm giờ trả thực tế. Luật hạn chót so giờ trả mới với lời hứa đầu, nên riêng phần thay đổi
khi giữ nguyên kế hoạch cũng có thể làm nó bị ép (ví dụ 40 giây ở mục~\ref{muc:doichung}).
% EV: A21; A22; N-21; X3-16 (U 0 / 20.760); X3-30; A10; K05

Tính gộp, luật hạn chót 30 giây bị ép ở 7.710 trong 20.103 lượt quyết định (38,4\%) và phá lời hứa với 289 trong 2.338
khách (12,4\%). Gánh nặng dồn vào khung 7 giờ: cả 9 kịch bản có khách bị phá, với trung bình theo kịch bản 36,5\%
khách bị phá và 75,6\% lượt quyết định bị ép. Ở khung 17 giờ chỉ 5 trong 11 kịch bản có khách bị phá, với 1,0\% và
5,0\% (cũng trung bình theo kịch bản). Ở mức 60 giây, các tỉ lệ toàn tập giảm nhưng không về 0. Luật trần độ dời khách
thấy bị ép nhiều nhất, phù hợp với định nghĩa: tổng độ lớn các lần cập nhật mà nó kiểm không nhỏ hơn độ lùi so với lời
hứa đầu (công thức~\eqref{eq:dongnhat}).
% EV: N-21 (gộp 7.710/20.103 = 38,35%; 289/2.338 = 12,36%; khung 7 giờ 9/9, 36,53%, 75,62%; 17 giờ 5/11, 1,04%, 4,99% — đều trung bình theo kịch bản); C14; C13 (M⁺-60 22,5% / 4,2%; V-30 53,8%, 34/34); A11; A12; X3-33 (eq:dongnhat: Σ|z| ≥ |Σz| = |p_n − p_0|)

Bảng~\ref{bang:c4-nghiavu} không so hai chuẩn đúng giờ: luật C hứa về phần dời do điều phối, còn luật hạn chót hứa về
toàn bộ độ trễ so với lời hứa đầu, một lời hứa lớn hơn. Đo bằng chính thước của luật hạn chót, tức tới muộn hơn lời
hứa đầu quá 30 giây, luật C có 292 trong 2.338 khách (12,5\%), so với 289 khách (12,4\%) của luật hạn chót. Trung bình
theo kịch bản là 14,6\% so với 14,5\%, và cả hai đều có khách như vậy ở 28 kịch bản. Vậy về độ đúng giờ, hai luật gần
như ngang nhau; luật C khác ở chỗ chỉ hứa phần mà bộ điều phối kiểm soát được, nên giữ được lời hứa đó.
% EV: C13; K05; C4A-01 (CK4A check: C-30 292/2.338 = 12,49%, TB theo kịch bản 14,59%, 28/34; khớp ADD1:17 late30 292 và TSV); N-21

Bảng~\ref{bang:c4-phanbotre} cho thấy số lượt trễ giảm ở cả bốn ngưỡng đã đo, từ 1 tới 10 phút.
% EV: N-24

\begin{table}[htbp]
\centering
\caption{Phân bố trễ so với lời hứa đầu không dự báo trên tập kiểm tra chính: số lượt trễ quá mỗi ngưỡng
(trong số khách được phục vụ) và số lượt đó tính trên 100 yêu cầu (3.672 yêu cầu mỗi cấu hình).}
\label{bang:c4-phanbotre}
\small
\begin{tabular}{@{}lrrrr@{}}
\toprule
& \multicolumn{2}{c}{\textbf{U}} & \multicolumn{2}{c}{\textbf{C-30}}\\
\cmidrule(lr){2-3}\cmidrule(lr){4-5}
\textbf{Trễ quá} & \textbf{Lượt} & \textbf{Trên 100 yêu cầu} & \textbf{Lượt} & \textbf{Trên 100 yêu cầu}\\
\midrule
1 phút & 203 & 5,53 & 76 & 2,07\\
2 phút & 102 & 2,78 & 11 & 0,30\\
5 phút & 18 & 0,49 & 0 & 0\\
10 phút & 2 & 0,05 & 0 & 0\\
\midrule
Trễ lớn nhất (s) & \multicolumn{2}{c}{660,1} & \multicolumn{2}{c}{232,3}\\
\bottomrule
\end{tabular}
\end{table}
% EV: N-24 (203/102/18/2 và 76/11/0/0; 5,53/2,78/0,49/0,05 và 2,07/0,30/0/0); C07; C28; C4A-14 (CON:4,11: +660,1 s và +232,3 s; khớp C26 "11 phút" và "3,87 phút")

Cách đọc bảng: khi không trần, 18 lượt tới muộn hơn lời hứa đầu quá 5 phút, tức 0,49 lượt trên 100 yêu cầu.
% EV: N-24; C28

Khi không trần, lượt trễ nhất muộn 660,1 giây (khoảng 11 phút); dưới trần 30 giây là 232,3 giây (3,87 phút). Trong 76
khách còn trễ quá một phút dưới luật C, 72 người không có lần dời do điều phối nào được ghi và 61 người thuộc khung 7
giờ. Lời hứa theo dự báo (mục~\ref{muc:c4-dubao}) nhắm tới phần đoán được của loại trễ này.
% EV: N-24; C4A-14; C26; C21 (4/76 từng bị dời); K39 (không viết "trễ thuần vì kẹt xe"); N-11 (trễ quá 60 s của C-30 theo khung 61 / 15 / 0)

Tách theo kịch bản, luật C không có nhiều lượt trễ quá một phút hơn điều phối không trần ở cả 34 kịch bản và có số
lượt không trễ quá một phút bằng hoặc nhiều hơn ở 29 kịch bản. Ở khung 8 và 17 giờ, điều phối không trần nhiều lượt
trễ hơn ở cả 25 kịch bản; ở khung 7 giờ, luật C ít hơn ở 3 kịch bản và bằng ở 6. Vì 34 kịch bản thuộc 5 ngày và không
độc lập, đây là phép đếm mô tả, không phải kiểm định thống kê.
% EV: N-22 (34/34 trễ ≤ U; 29/34 không trễ ≥ U; bảng theo ngày × khung ở ADD1:173-186); C30; I05; B11

% -----------------------------------------------------------------------------------------------------
\section{Phần do điều phối: luật C chặn gì và mất gì}
\label{muc:c4-dieuphoi}

Dưới luật C trần 30 giây, 33 trong 2.338 khách được phục vụ bị điều phối dời giờ trả, mỗi khách nhiều nhất 29,2 giây;
khi không trần là 174 trong 2.434 khách, một khách bị dời tổng cộng 590,1 giây (Bảng~\ref{bang:c4-doidieuphoi}). Cái
giá là số yêu cầu bị từ chối tăng từ 1.238 lên 1.334; chỉ 148 trong số đó có phương án bị trần chặn được ghi lại.
% EV: C20; N-23; C01 (1.238 = 3.672 − 2.434; CON:3); D13 (1.334; 148)

\begin{table}[htbp]
\centering
\caption{Phần dời giờ trả do điều phối trên tập kiểm tra chính, lời hứa đầu không dự báo. Mẫu số của
dòng đầu là số khách được phục vụ (2.434, 2.338 và 2.387); độ dời của một khách là tổng $|c|$ của khách
đó. Dòng cuối lấy tổng phần dời có dấu do điều phối chia cho tổng số giây trễ so với lời hứa đầu của chính các khách
trễ quá 60 s: 23.027 trên 31.353 giây, 82 trên 6.892 giây và 4.836 trên 12.450 giây. Cột trần 120 giây được dùng
lại ở mục~\ref{muc:c4-chontran}.}
\label{bang:c4-doidieuphoi}
\small
\begin{tabular}{@{}p{7.6cm}rrr@{}}
\toprule
\textbf{Chỉ số} & \textbf{U} & \textbf{C-30} & \textbf{C-120}\\
\midrule
Khách bị điều phối dời giờ trả & 174 (7,15\%) & 33 (1,41\%) & 103 (4,32\%)\\
Khách bị dời từ hai lần trở lên & 12 & 0 & 0\\
Tổng dời lớn nhất của một khách (s) & 590,1 & 29,2 & 118,9\\
Trung vị độ lớn mỗi lần dời (s) & 108,6 & 8,3 & 62,5\\
% EV: C20; C22 (trung vị theo lần: 108,6 s trên 187 lần, 8,3 s trên 33 lần); N-07 (C-120: 103, 0 khách ≥ 2 lần, nên trung vị theo lần = trung vị theo khách 62,54 s); N-23
Khách trễ quá 60 s từng bị dời & 134 trên 203 & 4 trên 76 & 60 trên 132\\
Phần tính cho điều phối trong tổng giây trễ của các khách trễ quá 60 s (\%) & 73,4 & 1,2 & 38,8\\
% EV: C21; C24; N-23 (cột C-120: 60/132, 38,8%); C4A-03 (tử số, mẫu số: 23.027,239 / 31.353,407; 81,742 / 6.891,526; 4.836,281 / 12.450,223 s)
\bottomrule
\end{tabular}
\end{table}

Cách đọc bảng: khi không trần, 174 khách bị điều phối dời giờ trả ít nhất một lần, 12 người trong số đó bị dời từ hai
lần, và tổng dời lớn nhất của một khách là 590,1 giây.
% EV: C20

Dưới trần 30 giây, mỗi khách bị dời chỉ một lần, trung vị 8,3 giây, so với 108,6 giây trên 187 lần dời khi không trần.
Trong 33 khách bị dời, 29 người bị dời muộn đi; 4 người bị kéo sớm hơn, mỗi người đúng 1 mili giây. Khi không trần,
134 trong 203 khách trễ quá một phút từng bị dời; con số chỉ cho biết hai nhóm trùng nhau, không cho biết việc bị dời
là nguyên nhân của độ trễ. Theo sổ theo từng khách (công thức~\eqref{eq:so}), phần tính cho điều phối bằng 73,4\% tổng
số giây trễ của 203 khách đó, còn dưới trần 30 giây chỉ 1,2\%: gần như mọi giây trễ còn lại nằm ở hai phần kia của sổ.
% EV: E02; C22; C23 (29 muộn, 4 sớm); N-05 (mỗi lần dời sớm đúng −1 ms; chỉ kiểm ở U, C-30, C-120, M⁺-30; nguyên nhân chưa kiểm, ADDENDUM-1 mục 0.2; C4A-19); C21 (giao hai tập); C24 (sổ mô hình, không nhân quả); K04; X3-12

Giới hạn này không dựa vào dự báo, vì dự báo không đi vào bộ điều phối. Ở 13 kịch bản cuối tuần không có dự báo, tổng
dời lớn nhất của một khách là 29,2 giây ở trần 30 giây và 116,5 giây ở trần 120 giây. Trên 8 kịch bản chiều tuyết (bộ
chạy chưa kiểm lại), 8 trong 601 khách được phục vụ bị dời, lớn nhất 25,2 giây; không ai vượt trần và không quyết định
nào bị ép trong 5.067 lượt.
% EV: F03; X3-11 (D chỉ tính sau khi chạy); N-28; N-26 (thay anchor F13 của RD §8 mục 4, ADDENDUM-1 mục 0.4); X3-10 (bộ chạy chiều tuyết chưa kiểm)

Một ví dụ trên tập buổi sáng cho thấy luật C và luật hạn chót khác nhau ở đâu. Trước một quyết định, giờ trả đang hứa
của một khách đã lùi 79,0 giây so với lời hứa đầu, toàn bộ là phần thay đổi khi giữ nguyên kế hoạch. Phương án ghép
thêm một khách mới sẽ dời giờ trả dự kiến của khách đó thêm 68,3 giây do điều phối. Ở mức 90 giây, luật hạn chót cộng
hai phần, khoảng 147 giây, nên loại phương án; luật C chỉ tính 68,3 giây nên cho qua. Trên tập này hai luật tách nhau
ở đúng ba chỗ cùng kiểu, ở mức 90 hoặc 120 giây; mỗi chỗ luật C nhận thêm một khách. Đây là phép tính trên giờ trả dự
kiến của mô hình, không phải phép đo nhân quả.
% EV: D05 (79,0 + 68,3 s; 60,9 + 61,2 s); C4A-12 (W07XM:6,11,16: "sum x 79.0, sum c 0.0, p_last − p_first 79.0", c = 68,3 s, x = 0,1 s; W07X:54-62: "khách đang có", giai đoạn chờ đón hay trên xe không ghi); K04

Khi luật C loại một phương án vì nó làm một khách đã có lời hứa vượt trần, bộ kiểm ghi một bản ghi chặn gồm khách đó,
phần trần đã dùng, lần dời mới và trần (Bảng~\ref{bang:kiemtran}). Dưới trần 30 giây, 148 trong 1.334 yêu cầu bị từ
chối có ít nhất một phương án mang bản ghi chặn; điều phối không trần phục vụ 116 trong số đó. 1.186 yêu cầu còn lại
không có bản ghi chặn nào được ghi, và điều phối không trần cũng từ chối 1.075 yêu cầu trong số này; điều đó không
chứng minh chúng bị từ chối chỉ vì ràng buộc dịch vụ. Bản ghi chặn cũng không phải khách: 845 bản ghi nằm ở 643 yêu
cầu, và 495 yêu cầu trong số đó cuối cùng vẫn được phục vụ.
% EV: X3-05 (nội dung bản ghi chặn); D13 (BUD:2-5: 148 = 116 + 32; 1.186 = 111 + 1.075); K10; C4A-09 (BUDPY:29-44: nhóm 1.186 là phần dư không có bản ghi chặn được ghi, không phải lý do cấu trúc đã chứng minh); D16 (SF:2: 643 cặp, 495 phục vụ, 148 từ chối); E01 (845)

Với mỗi yêu cầu trong 148 yêu cầu, phép phân tích chọn bản ghi có mức vượt nhỏ nhất, tức phần vượt quá 30 giây; trung
vị là 109,6 giây, lớn nhất 524,8 giây. Chỉ ở 2 trong 148 bản ghi được chọn, khách nêu trong bản ghi đã dùng một phần
trần trước đó; ở cả 148, riêng lần dời mới đã vượt 30 giây. Với trần 120 giây, mức vượt của 59 trong 148 yêu cầu nằm
trong phần nới thêm 90 giây. Không thể kết luận 59 yêu cầu đó sẽ được phục vụ: bộ kiểm dừng ở khách vi phạm gặp trước,
và khi trần đổi thì các quyết định trước cũng đổi. Mục~\ref{muc:c4-chontran} cho kết quả chạy thật.
% EV: D14 (BUD:6-8: trung vị 109,6 s; lớn nhất 524,8 s; 59/148 có mức vượt ≤ 90 s); D19 (BUD:10; FJC:4-6); C4A-09 (BUDPY:22,35-42,51-52: mức vượt = after − limit; 2/148 đếm theo bản ghi được chọn của từng yêu cầu, không phải số khách khác nhau); K11; A23; I08

% -----------------------------------------------------------------------------------------------------
\section{Chọn mức trần}
\label{muc:c4-chontran}

Trong ba điểm đã đo trên tập kiểm tra chính, không điểm nào tốt nhất ở mọi chỉ số. Trần 30 giây có ít lượt trễ quá một
phút nhất (76). Trần 120 giây có nhiều lượt không trễ quá hai phút nhất (2.366, so với 2.327 của trần 30 giây và 2.332
khi không trần) và phục vụ 2.387 lượt, nằm giữa hai điểm kia. Chọn mức nào tùy ngưỡng mà hãng xe coi trọng, điều thực
nghiệm không trả lời được.
% EV: N-06; L1a-02; C01; C06; C08; K12; K36

Mức trần được đo trên ba tập không cộng gộp (Hình~\ref{hinh:c4-tran}). Tập tắc đường cố định cho thấy trần làm gì khi
thời gian đi không đổi, tập buổi sáng dò tám mức trên cùng 8 kịch bản, tập kiểm tra chính thêm mức 120 giây.
% EV: B17; B15; B16; X3-28; D01; L1a-01; I12

\textbf{Tập tắc đường cố định.} Bảng~\ref{bang:c4-codinh} xếp các cấu hình từ lỏng nhất (ghép tự do) tới chặt nhất
(trần 30 giây), với hệ số tắc đường cố định suốt lần chạy.
% EV: B17; LAD:2-6

\begin{table}[htbp]
\centering
\caption{Các mức trần trên tập tắc đường cố định (16 kịch bản, 1.728 yêu cầu mỗi cấu hình; bộ chạy cũ,
không có luật phục hồi). Ba cột số đầu là số đếm; hai cột cuối là số khách được phục vụ thỏa điều kiện, trong ngoặc
là tỉ lệ gộp theo khách được phục vụ (\%).}
\label{bang:c4-codinh}
\small
\begin{tabular}{@{}lrrrrr@{}}
\toprule
\textbf{Cấu hình} & \textbf{Phục vụ} & \textbf{Trễ quá 60 s} & \textbf{Không trễ quá 120 s} &
\textbf{Đi chung} & \textbf{Bị dời giờ trả}\\
\midrule
FREE & 1.629 & 274 & 1.465 & 865 (53,1) & 388 (23,8)\\
U & 1.584 & 47 & 1.557 & 346 (21,8) & 62 (3,9)\\
C-120 & 1.563 & 23 & 1.563 & 303 (19,4) & 41 (2,6)\\
C-60 & 1.540 & 0 & 1.540 & 266 (17,3) & 20 (1,3)\\
C-30 & 1.534 & 0 & 1.534 & 241 (15,7) & 6 (0,4)\\
\bottomrule
\end{tabular}
\end{table}
% EV: D09 (số đếm, WP17); D08 (tỉ lệ đi chung, bị dời; gộp theo khách); C4A-04 (đi chung 865/346/303/266/241, CK4A fixed-shared; bị dời 388/62/41/20/6, WBU:4,7,10,13,16); B17

Cách đọc bảng: với trần 120 giây, 1.563 trong 1.728 yêu cầu được phục vụ, 23 lượt trễ quá 60 giây; 303 khách (19,4\%)
đi chung và 41 khách (2,6\%) bị dời giờ trả.
% EV: D08; D09; C4A-04

Ghép tự do khác điều phối không trần ở cả cách chọn phương án lẫn việc kiểm cam kết, nên chênh lệch giữa hai cấu hình
này không chỉ do trần. Trong nhóm cùng chính sách, siết từ không trần xuống trần 30 giây làm số lượt phục vụ giảm từ
1.584 xuống 1.534 và số lượt trễ quá 60 giây từ 47 xuống 0. Giữa hai mức đã chạy 60 và 120 giây, số lượt trễ tăng từ 0
lên 23. Các con số của bộ chạy cũ này không chuyển sang tập kiểm tra chính.
% EV: A17; K50; X3-38; NEW:13; D09; D11; B17; I12

\textbf{Tập chỉnh trần buổi sáng.} Trên 864 yêu cầu, số lượt phục vụ chỉ dao động trong 6 lượt, từ 660 ở trần 30 giây
tới 666 ở trần 120--180 giây (Bảng~\ref{bang:c4-sang}).
% EV: D01; B15; B16

\begin{table}[htbp]
\centering
\caption{Tám mức trần trên tập chỉnh trần buổi sáng (8 kịch bản khung 7 giờ ngày 12--13/11/2018, 864 yêu cầu
mỗi cấu hình; lời hứa đầu không dự báo). Số đếm gộp; hai cột ``không trễ quá'' là số lượt trên 864 yêu cầu. Các mức
có cùng mọi số đếm được gộp vào một dòng.}
\label{bang:c4-sang}
\small
\begin{tabular}{@{}lrrrr@{}}
\toprule
\textbf{Cấu hình} & \textbf{Phục vụ} & \textbf{Trễ quá 60 s} & \textbf{Không trễ quá 60 s} &
\textbf{Không trễ quá 120 s}\\
\midrule
U & 665 & 50 & 615 & 656\\
C-30 & 660 & 47 & 613 & 658\\
C-60 & 661 & 47 & 614 & 659\\
C-90 & 664 & 50 & 614 & 660\\
C-120, C-150, C-180 & 666 & 51 & 615 & 659\\
C-240, C-300 & 665 & 50 & 615 & 656\\
\bottomrule
\end{tabular}
\end{table}
% EV: D01; D02 (U, C-30, C-90, C-120: CAP:3-7; DIAL17:3-7, hai bộ đọc độc lập); C4A-05 (C-60, C-150, C-180, C-240, C-300: CAP:5,8-11; DIAL17:5,8-11); D02 lưu ý (cột dự báo bằng ngây thơ ở 12–13/11)

Cách đọc bảng: với trần 120, 150 hoặc 180 giây, 666 yêu cầu được phục vụ, 51 lượt trễ quá 60 giây, 615 lượt không trễ
quá 60 giây và 659 lượt không trễ quá 120 giây.
% EV: D01; D02; C4A-05

Trần 120 giây phục vụ nhiều hơn điều phối không trần một lượt, nên điều phối không trần không phải mức phục vụ cao
nhất có thể đạt, chỉ là một điểm đo. Từ 30 lên 120 giây, trần thêm 6 lượt phục vụ, 4 lượt trễ quá 60 giây và 2 lượt
không trễ quá 60 giây. Tỉ lệ lượt quyết định luật hạn chót bị ép (trung bình theo kịch bản) giảm từ 58,6\% ở mức 30
giây xuống 6,6\% ở 120 giây và 1,7\% ở 240--300 giây; luật C luôn 0\%. Ở cả tám mức, hai luật phục vụ chênh nhau không
quá 2 lượt.
% EV: D01; K12; D03 (+6 / +4 / +2); D04 (trung bình theo kịch bản; W07X:20-29); X3-28 (luật hạn chót chạy 30–300 s ở tập này)

\textbf{Tập kiểm tra chính.} Bảng~\ref{bang:c4-tran120} đặt trần 120 giây cạnh hai cấu hình của
mục~\ref{muc:c4-ketqua}, khi không có và khi có dự báo.

\begin{table}[htbp]
\centering
\caption{Trần 120 giây trên tập kiểm tra chính (34 kịch bản, 3.672 yêu cầu mỗi cấu hình). Số đếm gộp; các dòng
``không trễ quá'' là số lượt trên 3.672 yêu cầu. Khối dưới là phép phát lại lời hứa theo dự báo trên cùng các lần
chạy, bàn ở mục~\ref{muc:c4-dubao}; dự báo chỉ có ở 21 kịch bản ngày thường, và độ dài thêm của lời hứa là trung
bình gộp trên mọi khách được phục vụ, kể cả khách cuối tuần có độ dài thêm bằng 0.}
\label{bang:c4-tran120}
\small
\begin{tabular}{@{}p{8.4cm}rrr@{}}
\toprule
\textbf{Chỉ số} & \textbf{U} & \textbf{C-30} & \textbf{C-120}\\
\midrule
\multicolumn{4}{@{}l}{\textit{Lời hứa đầu không dự báo}}\\
Lượt phục vụ & 2.434 & 2.338 & 2.387\\
Lượt trễ quá 60 s & 203 & 76 & 132\\
Lượt trễ quá 120 s & 102 & 11 & 21\\
Lượt trễ quá 300 s & 18 & 0 & 0\\
Lượt tới sớm quá 60 s & 7 & 3 & 1\\
Không trễ quá 60 s & 2.231 & 2.262 & 2.255\\
Không trễ quá 120 s & 2.332 & 2.327 & 2.366\\
Khách đi chung & 555 & 313 & 428\\
% EV: L1a-02; L1a-04 (U, C-30 tái lập); N-06; C06; C07; C08; N-09 (L1a-07); C4A-13 (sớm quá 60 s 7 / 3 / 1, ADD1:12,17,22; C120R:14,16,18)
\midrule
\multicolumn{4}{@{}l}{\textit{Lời hứa đầu theo dự báo}}\\
Lượt trễ quá 60 s & 137 & 10 & 65\\
Lượt tới sớm quá 60 s & 8 & 4 & 2\\
Không trễ quá 60 s & 2.297 & 2.328 & 2.322\\
Không trễ quá 120 s & 2.347 & 2.338 & 2.379\\
Lời hứa dài thêm, trung bình gộp theo khách (s) & 7,18 & 7,37 & 7,27\\
% EV: L1a-03; N-03; N-06; N-16 (sớm quá 60 s 8 / 4 / 2); C09; F06
\bottomrule
\end{tabular}
\end{table}

Cách đọc bảng: với trần 120 giây và lời hứa đầu không dự báo, 2.387 yêu cầu được phục vụ, 132 lượt trễ quá 60 giây,
nên 2.255 lượt không trễ quá 60 giây trên 3.672 yêu cầu. Chỉ 1 lượt tới sớm quá 60 giây.
% EV: N-06; L1a-02; C4A-13

So với trần 30 giây, trần 120 giây phục vụ thêm 49 lượt và có thêm 56 lượt trễ quá một phút, nên ít hơn 7 lượt không
trễ quá một phút. Nhưng nó thêm 39 lượt không trễ quá hai phút, nhiều hơn cả trần 30 giây lẫn điều phối không trần.
Khi có dự báo, nó ít hơn trần 30 giây 6 lượt ở ngưỡng một phút và thêm 41 lượt ở ngưỡng hai phút. Độ dài thêm của lời
hứa gần như bằng nhau (7,27 so với 7,37 giây, trên hai tập khách khác nhau). Mục~\ref{muc:c4-dem} đặt điểm này cạnh
cách đệm lời hứa.
% EV: L1a-09 (+49 / +56 / −7 / +39; dự báo −6 / +41); N-06; L1a-02 (2.366 > U 2.332 và C-30 2.327); K12; L1a-03; N-16; L1a-08

Dưới trần 120 giây, 9 khách bị dời từ 108 giây trở lên (90\% của trần), không khách nào vượt trần và không quyết định
nào bị ép trong 20.435 lượt.
% EV: N-07 (≥ 0,9 B: 9); L1a-05; N-23

Hai tập khách khác nhau: 2.266 khách được cả hai mức trần phục vụ, 121 khách chỉ được trần 120 giây và 72 khách chỉ
được trần 30 giây phục vụ; 93 trong 121 khách đó cũng được phục vụ khi không trần. Vì vậy 49 lượt thêm là chênh lệch
hai tổng, không phải 49 trong 96 lượt mà trần 30 giây bỏ. Phần thêm dồn vào khung 8 và 17 giờ; so từng kịch bản, trần
120 giây phục vụ nhiều hơn ở 23, bằng ở 8 và ít hơn ở 3.
% EV: N-10 (lưu ý: không viết "lấy lại 49 trong 96 lượt"); L1a-08; N-11 (C-120 − C-30: phục vụ +4 / +26 / +19, trễ quá 60 s +2 / +34 / +20; 23 / 8 / 3 kịch bản); L1a-09

\begin{figure}[htbp]
\centering
\begin{tikzpicture}[font=\footnotesize,
  vuong/.style={draw,fill=white,rectangle,minimum size=4.2pt,inner sep=0pt},
  tamgiac/.style={draw,fill=white,regular polygon,regular polygon sides=3,minimum size=6pt,inner sep=0pt},
  tron/.style={draw,fill=white,circle,minimum size=4.6pt,inner sep=0pt},
  thoi/.style={draw,fill=black,diamond,minimum size=6.5pt,inner sep=0pt},
  nhan/.style={font=\scriptsize,inner sep=1.5pt}]
% (a) Tập tắc đường cố định. Trục ngang: phục vụ − 1.530, 0,06 cm mỗi lượt; trục dọc: trễ, 0,05 cm mỗi lượt.
\begin{scope}[shift={(0,0)}]
\draw[->] (0,0)--(3.95,0);
\draw[->] (0,0)--(0,3.35);
\foreach \vx/\so in {0/1.530,1.2/1.550,2.4/1.570,3.6/1.590}
  \draw (\vx,0)--(\vx,-0.07) node[below,nhan]{\so};
\foreach \vy/\so in {0/0,1/20,2/40,3/60}
  \draw (0,\vy)--(-0.07,\vy) node[left,nhan]{\so};
\draw[dashed] (0.24,0)--(0.60,0)--(1.98,1.15)--(3.24,2.35);
\node[vuong] at (0.24,0) {}; \node[nhan,above] at (0.24,0.05) {30};
\node[vuong] at (0.60,0) {}; \node[nhan,above] at (0.62,0.05) {60};
\node[vuong] at (1.98,1.15) {}; \node[nhan,above left] at (1.98,1.15) {120};
\node[vuong] at (3.24,2.35) {}; \node[nhan,above left] at (3.24,2.35) {U};
\node at (1.8,-0.72) {Lượt phục vụ};
\node[rotate=90] at (-0.85,1.6) {Lượt trễ quá 60 s};
\node[anchor=south] at (1.9,3.4) {(a) Tắc đường cố định};
\end{scope}
% (b) Tập chỉnh trần buổi sáng. Trục ngang: phục vụ − 658, 0,36 cm mỗi lượt; trục dọc: trễ, 0,05 cm mỗi lượt.
% C-120, C-150, C-180 trùng một điểm (666; 51); U, C-240, C-300 trùng một điểm (665; 50).
\begin{scope}[shift={(5.3,0)}]
\draw[->] (0,0)--(3.95,0);
\draw[->] (0,0)--(0,3.35);
\foreach \vx/\so in {0/658,0.72/660,1.44/662,2.16/664,2.88/666,3.6/668}
  \draw (\vx,0)--(\vx,-0.07) node[below,nhan]{\so};
\foreach \vy/\so in {0/0,1/20,2/40,3/60}
  \draw (0,\vy)--(-0.07,\vy) node[left,nhan]{\so};
\draw[densely dotted,thick] (0.72,2.35)--(1.08,2.35)--(2.16,2.50)--(2.88,2.55)--(2.52,2.50);
\node[tamgiac] at (0.72,2.35) {}; \node[nhan,below] at (0.72,2.25) {30};
\node[tamgiac] at (1.08,2.35) {}; \node[nhan,below] at (1.10,2.25) {60};
\node[tamgiac] at (2.16,2.50) {}; \node[nhan,below] at (2.16,2.40) {90};
\node[tamgiac] at (2.88,2.55) {}; \node[nhan,right] at (2.98,2.55) {120--180};
\node[tamgiac] at (2.52,2.50) {}; \node[nhan,above left] at (2.50,2.58) {U, 240, 300};
\node at (1.8,-0.72) {Lượt phục vụ};
\node[rotate=90] at (-0.85,1.6) {Lượt trễ quá 60 s};
\node[anchor=south] at (1.9,3.4) {(b) Buổi sáng};
\end{scope}
% (c) Tập kiểm tra chính. Trục ngang: phục vụ − 2.320, 0,03 cm mỗi lượt; trục dọc: trễ, 0,0125 cm mỗi lượt.
\begin{scope}[shift={(10.6,0)}]
\draw[->] (0,0)--(3.95,0);
\draw[->] (0,0)--(0,3.35);
\foreach \vx/\so in {0/2.320,1.2/2.360,2.4/2.400,3.6/2.440}
  \draw (\vx,0)--(\vx,-0.07) node[below,nhan]{\so};
\foreach \vy/\so in {0/0,1/80,2/160,3/240}
  \draw (0,\vy)--(-0.07,\vy) node[left,nhan]{\so};
\draw (0.54,0.95)--(2.01,1.65)--(3.42,2.54);
\node[tron] at (0.54,0.95) {}; \node[nhan,below right] at (0.54,0.95) {30};
\node[thoi] at (2.01,1.65) {}; \node[nhan,below right] at (2.06,1.62) {120};
\node[tron] at (3.42,2.54) {}; \node[nhan,above left] at (3.42,2.54) {U};
\node at (1.8,-0.72) {Lượt phục vụ};
\node[rotate=90] at (-0.85,1.6) {Lượt trễ quá 60 s};
\node[anchor=south] at (1.9,3.4) {(c) Tập kiểm tra chính};
\end{scope}
\end{tikzpicture}
\caption{Lượt trễ quá 60 giây so với lời hứa đầu (trục dọc) theo lượt phục vụ (trục ngang) ở các mức
trần, lời hứa đầu không dự báo. Mỗi khung là một tập với thang riêng: (a) tập tắc đường cố định, 1.728
yêu cầu; (b) tập chỉnh trần buổi sáng, 864 yêu cầu, đủ tám mức, các mức có cùng số đếm chung một điểm; (c) tập
kiểm tra chính, 3.672 yêu cầu. Số cạnh mỗi điểm là trần tính bằng giây, U là điều phối không trần. Đường nối các mức
của cùng một tập theo thứ tự trần tăng dần, chỉ giúp theo dõi các điểm đã đo và không biểu diễn kết quả ở các mức
chưa chạy; các tập không nối với nhau. Ở (a), ghép tự do (1.629 lượt phục vụ, 274 lượt trễ) nằm ngoài khung. Hình
thoi đen là trần 120 giây trên tập kiểm tra chính.}
\label{hinh:c4-tran}
\end{figure}
% EV (a): D09 (C-30 1.534/0; C-60 1.540/0; C-120 1.563/23; U 1.584/47; FREE 1.629/274)
% EV (b): D01; D02; C4A-05 (C-30 660/47; C-60 661/47; C-90 664/50; C-120 = C-150 = C-180 666/51; U = C-240 = C-300 665/50)
% EV (c): C01; C06; N-06 (C-30 2.338/76; C-120 2.387/132; U 2.434/203)
% Tọa độ: (a) x = (phục vụ − 1530) × 0,06, y = trễ × 0,05; (b) x = (phục vụ − 658) × 0,36, y = trễ × 0,05
% (C-60: x = 3 × 0,36 = 1,08; y = 47 × 0,05 = 2,35); (c) x = (phục vụ − 2320) × 0,03, y = trễ × 0,0125.

Trên cả ba tập, nới trần từ 30 lên 120 giây thêm lượt phục vụ và thêm lượt trễ quá một phút. Số lượt không trễ quá một
phút thêm 6, thêm 2 và bớt 7 lần lượt ở tập tắc đường cố định, tập buổi sáng và tập chính, còn ở ngưỡng hai phút trần
120 giây không kém trần 30 giây ở cả ba tập. Vì mỗi kịch bản chỉ chạy một lần và mức 120 giây được chọn sau khi đã xem
tập kiểm tra chính, trần 120 giây chỉ là một điểm đo thêm trên đường đánh đổi.
% EV: D10; D03; L1a-09; D09 (không trễ quá 120 s: 1.563 so với 1.534); D02 (659 so với 658); L1a-02; K36; D11; B11; I06; C4A-11; K42

% -----------------------------------------------------------------------------------------------------
\section{Phần cộng dồn của luật C: kiểm toàn bộ bản ghi chặn}
\label{muc:c4-congdon}

Luật C cộng mọi lần dời của một khách qua cả chuyến rồi mới so với trần. Trên các lần chạy đã kiểm, hiếm khi chỉ riêng
phép cộng ấy làm một phương án bị loại. Trong 98 lần chạy trần 30--300 giây của tập kiểm tra chính và tập buổi sáng,
cả 1.256 bản ghi chặn đều có một lần dời đơn lẻ đã vượt trần của cấu hình tương ứng. Trên 34 lần chạy trần 120 giây
của tập kiểm tra chính, 3 trong 524 bản ghi chỉ vượt trần khi cộng thêm phần trần đã dùng trước đó.
% EV: E01 (98 = 34 + 64; 1.256 = 845 + 411); K09; N-08; L1a-06

Một bản ghi chặn là \emph{chỉ do cộng dồn} khi lần dời mới không vượt $B$ nhưng phần đã dùng $L_r$ cộng lần dời mới
thì vượt (mục~\ref{muc:luatc}). Ví dụ, ở một bản ghi ngày 14/11 với trần 120 giây, khách đã dùng 25,844 giây và phương
án mới dời thêm 99,027 giây: mỗi phần dưới 120 giây, nhưng tổng 124,871 giây vượt trần. Bảng~\ref{bang:c4-congdon} đếm
các bản ghi này trên mọi bản ghi chặn của các lần chạy đã kiểm.
% EV: A06; E05; N-27 (25,844 + 99,027 = 124,871 s); N-08 (quy tắc riêng + cumulative_decisive_check)

\begin{table}[htbp]
\centering
\caption{Kiểm toàn bộ bản ghi chặn do trần trên 140 lần chạy có luật C dùng hệ số tắc đường đo được đã được kiểm,
thuộc bốn nhóm dưới đây. Một bản ghi ứng với một phương án bị loại, không phải một khách.}
\label{bang:c4-congdon}
\small
\begin{tabular}{@{}p{6.0cm}rrrr@{}}
\toprule
\textbf{Lần chạy} & \textbf{Số lần chạy} & \textbf{Bản ghi chặn} & \textbf{Chỉ do cộng dồn} &
\textbf{\begin{tabular}[b]{@{}r@{}}Khách phục vụ bị\\ dời từ hai lần\end{tabular}}\\
\midrule
Tập kiểm tra chính, trần 30 giây & 34 & 845 & 0 & 0\\
Tập chỉnh trần buổi sáng, trần 30--300 giây & 64 & 411 & 0 & 0\\
% EV: E01 (98 = 34 + 64; 1.256 = 845 + 411; 0 chỉ do cộng dồn); E02 (repeated = 0 trên 98 lần chạy)
Tập kiểm tra chính, trần 120 giây & 34 & 524 & 3 & 0\\
% EV: N-08; L1a-06; N-07 (0 khách ≥ 2 lần)
Chiều tuyết 15/11/2018, trần 30 giây & 8 & 163 & 0 & 0\\
% EV: N-26 (163 bản ghi, 0 chỉ do cộng dồn, ≥ 2 lần 0); X3-10 (bộ chạy chưa kiểm)
\midrule
Cộng & 140 & 1.943 & 3 & 0\\
% EV: N-15 (132 / 1.780 / 3; kể cả chiều tuyết 140 / 1.943 / 3); ADDENDUM-1 mục 0.3
\bottomrule
\end{tabular}
\end{table}

Cách đọc bảng: trên 34 lần chạy trần 120 giây có 524 phương án bị loại kèm bản ghi chặn do trần, 3 bản ghi trong số đó
chỉ do cộng dồn, và không khách được phục vụ nào bị dời từ hai lần.
% EV: N-08; N-07

Phạm vi kiểm là 132 lần chạy và 1.780 bản ghi (98 lần chạy với 64.371 lượt quyết định được xác nhận, cộng 34 lần chạy
trần 120 giây), hay 140 lần chạy và 1.943 bản ghi kể cả chiều tuyết. Chưa được kiểm: 15 lần chạy hiệu chỉnh, 16 lần
chạy trong đó bộ điều phối dùng dự báo khi quyết định (khác phép phát lại) và 12 lần chạy luật C dùng hệ số đo được
của một phép quét trước đó. Kết quả vì vậy không áp cho mọi lần chạy.
% EV: N-15 (SCOPE); ADDENDUM-1 mục 0.3 (15 lần chạy hiệu chỉnh, 16 thư mục C-30F); E01 (64.371); K09; F12; C4A-08 (SWEEP:13,70: C-x, x = 30…900 s, trên thế giới đo thật W5, 2 kịch bản mỗi mức; NEW:1 by_arm chỉ có t38 và w07)

Ba bản ghi chỉ do cộng dồn nằm ở ba kịch bản khung 8 giờ (14, 15 và 18/11), trong 3 lượt quyết định; khi không trần,
khách của bản ghi ngày 14/11 bị dời hai lần đúng bằng hai lượng ở ví dụ trên. Trong 3 yêu cầu mới liên quan, 2 cuối
cùng không được nhận dưới trần 120 giây (một trong hai cũng không dưới trần 30 giây). Vì bản ghi dừng ở lỗi gặp trước
và quỹ đạo các cấu hình khác nhau, điều này không chứng minh phần cộng dồn làm mất hai yêu cầu đó.
% EV: L1a-06 (ba kịch bản w08 ngày 14, 15, 18/11); N-08 (3 quyết định; 2 không được nhận, 1 phục vụ); N-27; E07; A23; I08

Ngoài hệ số đo được, bản ghi chỉ do cộng dồn còn có ở hai nơi. Trên tập tắc đường cố định, nơi bộ kiểm ghi đủ mọi bản
ghi, có 2 bản ghi ở trần 60 giây và 4 ở trần 120 giây; mọi yêu cầu mới liên quan vẫn được phục vụ. Trên các kiểu tắc
đường tự tạo của phép quét trước đó, sau khi bỏ bản chạy trùng, có 4 bản ghi trong 3 lượt quyết định ở trần 120 giây;
1 yêu cầu mới liên quan cuối cùng bị từ chối.
% EV: E03; E04 (4/3 sau khi bỏ bản lặp; thô 6/5); A23 (NEW:13: tập tắc đường cố định bật chế độ ghi đủ); FJC:31-54; FJR:2-7; C4A-08 (thế giới tự tạo thuộc tier3-sweep-v1)

Về định nghĩa, luật C vẫn khác luật ``mỗi lần dời không quá 30 giây'': luật sau cho một khách bị dời ba lần, tổng 90
giây, còn luật C trần 30 giây thì không. Trên số liệu, khác biệt này hiếm: khi không trần chỉ 12 trong 2.434 khách bị
dời từ hai lần. Tính thử cả hai cách trên các quỹ đạo không trần đã ghi, chúng chỉ khác kết luận ở 4 khách với trần
120 giây và 0 khách với trần 30, 60 và 300 giây. Phép tính thử không cho biết kết quả nếu chạy thật với luật kia, vì
các quyết định sẽ khác; luật mỗi lần một trần chưa được cài đặt và chưa chạy.
% EV: E08; C20; E06 (CVP:14-18: 0/0/4/0 trên 2.434 khách U); K42; E09; K08

Các mục trên so luật C với điều phối không trần khi giữ nguyên bộ chọn ưu tiên ổn định, tức bộ chọn xét mức dùng trần
và phần dời trước chi phí kế hoạch (Bảng~\ref{bang:rangbuoc}). Mục~\ref{muc:c4-vaitro} xét kết quả có phụ thuộc vào
thứ tự đó không.
% EV: A08; H01; H02; A14; X3-07
